\documentclass[11pt]{article}
\usepackage[margin=2.1cm]{geometry}
\usepackage{booktabs}
\usepackage{amsmath}
\usepackage{amssymb}
\usepackage{enumitem}
\usepackage{array}
\usepackage{listings}
\usepackage{xcolor}
\usepackage[T1]{fontenc}
\usepackage{lmodern}

\setlength{\parindent}{0pt}
\setlength{\parskip}{3.6pt}
\newcommand{\code}[1]{\texttt{#1}}
\newcommand{\hr}{\vspace{2pt}\hrule\vspace{6pt}}
\renewcommand{\arraystretch}{1.12}

\definecolor{kw}{RGB}{0,90,160}
\definecolor{cmt}{RGB}{110,110,110}
\lstdefinestyle{vhdl}{
  language=VHDL, basicstyle=\ttfamily\footnotesize,
  keywordstyle=\color{kw}\bfseries, commentstyle=\color{cmt}\itshape,
  showstringspaces=false, columns=fullflexible, breaklines=true,
  frame=leftline, framesep=6pt, xleftmargin=8pt, aboveskip=6pt, belowskip=6pt
}
\lstset{style=vhdl}

\begin{document}

\begin{center}
{\Large\bfseries EE\,362 --- Pre-Lab 5: Solution}\\[2pt]
{\large Arithmetic Logic Unit}\\[4pt]
Fall 2026 \quad$\cdot$\quad 5 points \quad$\cdot$\quad TA marking guide
\end{center}
\hr

{\small Marking: Q1 = 2 points, Q2 = 1 point, Q3 = 2 points. The ALU is the
biggest component so far; students who arrive with Q1 and Q3 done finish the
session.}

\vspace{4pt}\hr

\section*{Question 1 --- Operations by hand (2 points)}

\subsection*{(a) The six results \quad(\code{Shft} = 1)}

\begin{center}
\begin{tabular}{clll|l|l}
\toprule
\# & op & \code{Ain} & \code{Bin} & \code{Rout} & why \\
\midrule
1 & \code{SUB}  & \code{0x00000004} & \code{0x00000007} & \code{0xFFFFFFFD} & $4-7 = -3$ \\
2 & \code{NOT}  & \code{0x00000000} & \code{0x0F0F0F0F} & \code{0xF0F0F0F0} & \code{NOT} acts on \code{Bin} \\
3 & \code{SRL}  & \code{0xFFFFFFF8} & --- & \code{0x7FFFFFFC} & vacated bit filled with \code{0} \\
4 & \code{SRA}  & \code{0xFFFFFFF8} & --- & \code{0xFFFFFFFC} & vacated bit filled with the sign; $-8/2 = -4$ \\
5 & \code{SOLS} & \code{0xFFFFFFFF} & \code{0x00000001} & \code{0x00000001} & signed: $-1 < 1$ is true \\
6 & \code{SOLU} & \code{0xFFFFFFFF} & \code{0x00000001} & \code{0x00000000} & unsigned: $4294967295 < 1$ is false \\
\bottomrule
\end{tabular}
\end{center}

\textbf{Row 2 is a trap worth marking carefully.} In this ALU \code{NOT} operates
on \code{Bin}, not \code{Ain}. A student who answered \code{0xFFFFFFFF} inverted
\code{Ain} instead --- correct arithmetic, wrong operand, and it will fail the
reference testbench.

\textbf{Rows 3 and 4 are the heart of the lab.} If a student wrote the same answer
twice, they have not yet seen the difference between the two right shifts.

\subsection*{(b) Same inputs, different answers}

\begin{quote}
\textbf{Rows 3 vs 4:} whether the bits are a \textbf{signed} number or an
\textbf{unsigned} one. \code{SRA} must preserve the sign, so it fills with copies
of the top bit; \code{SRL} treats the bits as unsigned and fills with zeros.
\end{quote}

\begin{quote}
\textbf{Rows 5 vs 6:} the same distinction applied to the comparison. As a signed
number \code{0xFFFFFFFF} is $-1$, which is less than 1. As an unsigned number it
is the largest 32-bit value, which is not.
\end{quote}

\textbf{The general point}, which is the same as Lab 4's: \emph{the bits do not
say what they mean; the operation does.}

\subsection*{(c) The two flag equations}

\[
\text{Zero} = \overline{R_{31} + R_{30} + \cdots + R_1 + R_0}
\qquad\qquad
\text{Neg} = R_{31}
\]

Zero is a 32-input NOR: it is 1 only when every bit is 0. Accept
$\text{Zero} = \overline{R_{31}} \cdot \overline{R_{30}} \cdots \overline{R_0}$ ---
De Morgan's law makes them the same thing, and a student who writes it that way
should be able to say so.

\begin{quote}
\textbf{Why \code{Neg} needs no computation:} in two's complement the top bit
\emph{is} the sign. \code{Neg} is a wire from $R_{31}$, not a gate.
\end{quote}

\vspace{4pt}\hr

\section*{Question 2 --- Interface and the flag trap (1 point)}

\subsection*{(a) The entity}

\begin{lstlisting}
entity ALU is
    port (
        Ain  : in  std_logic_vector(31 downto 0);
        Bin  : in  std_logic_vector(31 downto 0);
        Func : in  std_logic_vector(4 downto 0);   -- 16 operations, 5 bits
        Shft : in  std_logic_vector(4 downto 0);   -- shift amount 0..31
        Rout : out std_logic_vector(31 downto 0);
        Zero : out std_logic;
        Neg  : out std_logic
    );
end entity ALU;
\end{lstlisting}

\textbf{\code{Shft} is a separate port, and shifts act on \code{Ain}.} Students
who assume the shift amount comes from \code{Bin} will pass their own testbench
and fail in Lab 12 --- worth correcting now, in ink, on the pre-lab.

Five bits for \code{Shft} is exactly $0 \ldots 31$: nothing wasted, nothing out of
range.

\subsection*{(b) Taking \code{Zero} from the adder instead of the selected result}

\begin{quote}
\textbf{Model answer:} \code{AND} with \code{Ain = 0x0F0F0F0F} and
\code{Bin = 0xF0F0F0F0}. The selected result is \code{0x00000000}, so \code{Zero}
\emph{should} be 1. The adder computes \code{0xFFFFFFFF}, which is not zero, so
the flag would read 0 --- wrong.
\end{quote}

Any operation whose result differs from the sum earns full credit. The converse
example is just as good: \code{AND} of \code{0x00000001} and \code{0xFFFFFFFF}
gives 1 (\code{Zero} should be 0), while the adder's sum is \code{0x00000000} and
would report \code{Zero} = 1.

\begin{quote}
\textbf{The rule:} \textbf{flags describe the selected result}, so they are
computed \emph{after} the multiplexer, never inside one branch.
\end{quote}

In Lab 8 six one-line Boolean equations over \code{Zero} and \code{Neg} produce
every conditional branch the machine has. A flag taken from the wrong place there
breaks branching, not arithmetic --- which is much harder to find.

\vspace{4pt}\hr

\section*{Question 3 --- Test plan (2 points)}

\subsection*{The coverage argument}

\begin{quote}
\textbf{Model answer:} $2^{64}$ operand pairs cannot be enumerated, so test
\textbf{every operation at least once} --- the \code{case} is a multiplexer, and
an untested branch is an unverified input to it --- and for each operation use
\textbf{boundary operands} rather than random ones, because arithmetic bugs live
at carries, at the sign bit and at overflow. Then check the flags on a result
that is \emph{not} the adder's.
\end{quote}

\subsection*{Model vector list}

\begin{center}
\small
\begin{tabular}{c>{\raggedright\arraybackslash}p{5.0cm}>{\raggedright\arraybackslash}p{7.6cm}}
\toprule
\# & operation and operands & which wrong implementation it catches \\
\midrule
1 & \code{ADD}, \code{0x7FFFFFFF} + \code{0x00000001} & a carry that is dropped at the sign bit \\
2 & \code{SUB}, $4 - 7$ & operands swapped: gives $+3$ instead of \code{0xFFFFFFFD} \\
3 & \code{SRA}, \code{0xFFFFFFF8}, \code{Shft}=1 & \textbf{\code{SRA} written with \code{unsigned()}} --- it silently becomes \code{SRL} \\
4 & \code{SRL}, same operand and amount & the reverse swap, and it pins both shifts at once \\
5 & \code{SOLS} vs \code{SOLU}, \code{0xFFFFFFFF} vs \code{0x00000001} & a signed comparison used where unsigned was meant \\
6 & \code{AND}, \code{0x0F0F0F0F} \& \code{0xF0F0F0F0}, check \code{Zero} & \textbf{flags taken from the adder} instead of the selected result \\
7 & shift with \code{Shft} = 0 and \code{Shft} = 31 & an off-by-one in the shift amount conversion \\
\bottomrule
\end{tabular}
\end{center}

Six vectors earn full credit. \textbf{Vectors 3 and 6 are the ones that must
appear.}

\subsection*{(a) Which vector catches an \code{SRA} written as an \code{SRL}?}

\begin{quote}
\textbf{Model answer:} Vector 3 --- \code{SRA} of a \textbf{negative} operand with
a non-zero shift. \code{0xFFFFFFF8} shifted right by one gives \code{0xFFFFFFFC}
($-4$) if the fill is the sign bit, and \code{0x7FFFFFFC} if it is zeros.
\end{quote}

\textbf{A positive operand cannot catch it}, because both shifts fill with zeros
when the sign bit is already 0. Neither can a shift amount of zero. If a student's
only shift vector uses a positive number, the bug survives their entire testbench
--- and in VHDL it is an easy bug to write, since \code{shift\_right} is the same
operator for both and only the \emph{type} decides the fill.

\subsection*{(b) The five boundary values, and why the middle is dull}

\begin{center}
\begin{tabular}{ll}
\toprule
value & why it is a boundary \\
\midrule
\code{0x00000000} & zero: the only value that sets the \code{Zero} flag \\
\code{0x00000001} & smallest positive; catches off-by-one \\
\code{0x7FFFFFFF} & largest positive; $+1$ here overflows into the sign bit \\
\code{0x80000000} & most negative; its negation cannot be represented \\
\code{0xFFFFFFFF} & $-1$ signed, maximum unsigned: separates \code{SOLS} from \code{SOLU} \\
\bottomrule
\end{tabular}
\end{center}

\begin{quote}
\textbf{Why the middle of the range proves little:} every value in the middle
exercises the same carry chain in the same way. Bugs live where behaviour
\emph{changes} --- at a carry out of the top bit, at the transition from positive
to negative, and at the extremes of the shift amount. Testing \code{0x12345678}
plus \code{0x11111111} tells you almost nothing that \code{1 + 1} did not.
\end{quote}

\vspace{6pt}
\hr
{\small \textbf{At the bench:} the two answers that predict trouble are a Q1 with
rows 3 and 4 identical, and a Q3 with no negative-operand shift. Both students
will produce an ALU that passes their own tests and fails in Lab 12.}

\end{document}
